% Options for packages loaded elsewhere
\PassOptionsToPackage{unicode}{hyperref}
\PassOptionsToPackage{hyphens}{url}
\documentclass[
]{article}
\usepackage{xcolor}
\usepackage{amsmath,amssymb}
\setcounter{secnumdepth}{-\maxdimen} % remove section numbering
\usepackage{iftex}
\ifPDFTeX
  \usepackage[T1]{fontenc}
  \usepackage[utf8]{inputenc}
  \usepackage{textcomp} % provide euro and other symbols
\else % if luatex or xetex
  \usepackage{unicode-math} % this also loads fontspec
  \defaultfontfeatures{Scale=MatchLowercase}
  \defaultfontfeatures[\rmfamily]{Ligatures=TeX,Scale=1}
\fi
\usepackage{lmodern}
\ifPDFTeX\else
  % xetex/luatex font selection
\fi
% Use upquote if available, for straight quotes in verbatim environments
\IfFileExists{upquote.sty}{\usepackage{upquote}}{}
\IfFileExists{microtype.sty}{% use microtype if available
  \usepackage[]{microtype}
  \UseMicrotypeSet[protrusion]{basicmath} % disable protrusion for tt fonts
}{}
\makeatletter
\@ifundefined{KOMAClassName}{% if non-KOMA class
  \IfFileExists{parskip.sty}{%
    \usepackage{parskip}
  }{% else
    \setlength{\parindent}{0pt}
    \setlength{\parskip}{6pt plus 2pt minus 1pt}}
}{% if KOMA class
  \KOMAoptions{parskip=half}}
\makeatother
\usepackage{longtable,booktabs,array}
\newcounter{none} % for unnumbered tables
\usepackage{calc} % for calculating minipage widths
% Correct order of tables after \paragraph or \subparagraph
\usepackage{etoolbox}
\makeatletter
\patchcmd\longtable{\par}{\if@noskipsec\mbox{}\fi\par}{}{}
\makeatother
% Allow footnotes in longtable head/foot
\IfFileExists{footnotehyper.sty}{\usepackage{footnotehyper}}{\usepackage{footnote}}
\makesavenoteenv{longtable}
\setlength{\emergencystretch}{3em} % prevent overfull lines
\providecommand{\tightlist}{%
  \setlength{\itemsep}{0pt}\setlength{\parskip}{0pt}}
\usepackage{bookmark}
\IfFileExists{xurl.sty}{\usepackage{xurl}}{} % add URL line breaks if available
\urlstyle{same}
\hypersetup{
  hidelinks,
  pdfcreator={LaTeX via pandoc}}

\author{}
\date{}

\begin{document}

\section{Title Page (separate file)}\label{title-page-separate-file}

\section{(Submit as a SEPARATE file; do NOT include in the anonymized
manuscript.)}\label{submit-as-a-separate-file-do-not-include-in-the-anonymized-manuscript.}

\textbf{Title.} Compound-Truck Cross-Docking Scheduling with Arrival
Time Windows: A Bottleneck-Guided Iterated Local Search with
Exact-Solver Verification

\textbf{Authors.} Youngsoo Kim ᵃ, Sangjin Kwon ᵃ,*

\textbf{Affiliations.} ᵃ Department of Artificial Intelligence, Yonsei
University, 50 Yonsei-ro, Seodaemun-gu, Seoul 03722, Republic of Korea.

\textbf{Corresponding author (*).} Sangjin Kwon, Department of
Artificial Intelligence, Yonsei University, Seoul 03722, Republic of
Korea. Email: {[}advisor email --- to be added{]}

\textbf{Acknowledgements.} {[}Optional --- e.g., computing resources or
non-author help.{]}

\textbf{CRediT author statement.} - \textbf{Youngsoo Kim:}
Conceptualization, Methodology, Software, Formal analysis,
Investigation, Data curation, Visualization, Writing -- original draft.
- \textbf{Sangjin Kwon:}

\textbf{Declaration of competing interests.} The authors declare that
they have no known competing financial interests or personal
relationships that could have appeared to influence the work reported in
this paper.

\textbf{Funding.} This research did not receive any specific grant from
funding agencies in the public, commercial, or not-for-profit sectors.

\textbf{Data availability.} The instance generator, solver code, and
result files that support the findings of this study are available in a
public repository {[}URL to be added at acceptance{]}.

\newpage

\section{Highlights (separate file)}\label{highlights-separate-file}

\section{(Submit as a SEPARATE file; 3--5 bullets, each ≤ 85
characters.)}\label{submit-as-a-separate-file-35-bullets-each-85-characters.}

\begin{itemize}
\tightlist
\item
  First compound-truck cross-dock model with release times and soft due
  dates
\item
  MILP and CP-SAT formulations verify the heuristic on small and medium
  cases
\item
  VAA-GILS matches CP-SAT incumbents within 0.6\% in seconds, not
  minutes
\item
  Best-improvement descent and bottleneck moves drive the quality gains
\item
  Learned operator selection adds no benefit over uniform under equal
  budget
\end{itemize}

\newpage

\section{Manuscript (anonymized)}\label{manuscript-anonymized}

\textbf{Title.} Compound-Truck Cross-Docking Scheduling with Arrival
Time Windows: A Bottleneck-Guided Iterated Local Search with
Exact-Solver Verification

\subsection{Abstract}\label{abstract}

This paper considers a multi-door cross-dock in which a truck can be
partially unloaded and then leave as an outbound carrier. Existing
compound-truck models place every truck at the dock at time zero and
optimize makespan. Here, trucks have individual release times and soft
departure due dates, and the objective also charges tardiness. We
formulate the problem as a mixed-integer program and as a CP-SAT model,
and use separate train, tuning, and test seeds in the computational
study. Our method, VAA-GILS, starts from a Vogel-style construction and
uses best-improvement descent throughout; its additional moves focus on
the door that determines makespan or on the most tardy truck, with
simulated-annealing acceptance and kick restarts for diversification. On
the subsets for which CP-SAT produced a reference, the resulting
objectives were 0.1--0.6\% from the CP-SAT incumbents, and the small
no-window references were proven optimal. A GILS run took 0.1--2.3
seconds, versus 76--376 seconds for CP-SAT. For the larger time-window
cases, CP-SAT found no feasible schedule within 600 seconds, while GILS
returned the best solutions among the methods tested. Ablation is less
favorable to learning: descent, the bottleneck moves, and kick restarts
account for the improvement, whereas changing the initial solution,
retaining stochastic acceptance, or replacing uniform operator sampling
with tabular Q-learning or a transfer DQN changed the objective by no
more than 0.2 percentage points, and the learned selectors did not win
at any tested budget.

\textbf{Keywords:} cross-docking; truck scheduling; time windows;
iterated local search; constraint programming; guided local search

\subsection{1. Introduction}\label{introduction}

Cross-docking centers transfer products from inbound to outbound trucks
with little or no storage, reducing inventory and lead time (Van Belle
et al., 2012). In the compound-truck variant introduced by Joo and Kim
(2013) and refined by Shahmardan and Sajadieh (2020), a truck may serve
both roles: a compound truck arrives loaded with products for several
destinations, is \emph{partially} unloaded --- retaining the demand of
one destination --- and then reloads that destination's remaining demand
before departing as an outbound truck. Partial unloading was shown to
reduce makespan by up to 56\% relative to full unloading (Shahmardan \&
Sajadieh, 2020), making this model practically attractive for
less-than-truckload consolidation networks.

The simultaneous-arrival assumption is restrictive for scheduled docks,
where release times differ by truck. A makespan-only objective is also
incomplete when outbound departures have due dates. Release times, time
windows, and due-date objectives have been studied for conventional
inbound/outbound trucks (Assadi \& Bagheri, 2016; Bodnar et al., 2017;
Konur \& Golias, 2013a, 2013b; Ladier \& Alpan, 2016b; Larbi et al.,
2011; Molavi et al., 2018; Van Belle et al., 2013; Xi et al., 2020); our
contribution is to integrate known release times and soft due dates into
the partial-unloading compound-truck setting. To our knowledge, this
combination has not previously been formulated. Section 3 gives the
big-M model, its reified CP-SAT counterpart, and lower bounds. The
instance generator and the train/tuning/test seed split are included so
that the computational protocol can be reproduced.

For this problem we develop VAA-GILS. The method keeps the VAA
construction and generic moves used in earlier work (Shahmardan \&
Sajadieh, 2020), but adds a full best-improvement descent,
bottleneck-specific moves, and restarts. Its default selector is uniform
rather than learned. The exact-solver comparison is deliberately limited
to the subsets on which CP-SAT returned an incumbent; only the small
no-window subset has proven optima. Elsewhere we compare against the
best solutions found by the tested methods.

Much of the paper is devoted to identifying which parts of GILS matter.
The operator-pool experiment gives a monotone improvement when critical
and tardy moves are added (p \textless{} 0.002). Removing descent causes
the largest loss, followed by removing kick restarts. In contrast, the
final quality changes little when VAA is replaced by a random feasible
start or when simulated-annealing acceptance is removed. We obtained a
similar result for operator selection: neither tabular Q-learning nor a
transfer-trained DQN gave a meaningful gain over uniform sampling over
budgets of 50--3,000 iterations.

These comparisons concern our reimplementation and benchmark. They
should not be read as a replication claim about Shahmardan and Sajadieh
(2020), whose experimental setting and code differ. Their narrower
implication is methodological: when a learned selector is embedded in a
strong search engine, a uniform, equal-budget control is needed to
establish its incremental value.

\subsection{2. Related work}\label{related-work}

\textbf{Cross-dock truck scheduling with temporal constraints.} Boysen
and Fliedner (2010) classify dock-scheduling problems; Van Belle et
al.~(2012) and Ladier and Alpan (2016a) survey the field and its
industry gap. Several deterministic models already incorporate temporal
constraints for conventional inbound and outbound trucks. Van Belle et
al.~(2013) consider multiple docks, time windows, and tardiness; Assadi
and Bagheri (2016) use truck ready times and minimize outbound earliness
and tardiness; and Bodnar et al.~(2017) combine time windows with
mixed-service doors. Molavi et al.~(2018) instead impose hard outbound
due dates, while Rijal et al.~(2019) integrate truck scheduling and door
assignment and include tardiness in a mixed-service-door setting. These
studies establish that releases and due-date performance are not new to
cross-dock scheduling. They retain the conventional separation between
inbound and outbound trucks, however, and do not model a partially
unloaded truck that subsequently becomes a carrier.

\textbf{Incomplete and uncertain arrival information.} Arrival
uncertainty is a separate, well-developed stream. Konur and Golias study
both bounded arrival windows (2013a) and cost-stable schedules under
unknown arrivals (2013b). Larbi et al.~(2011) distinguish full, partial,
and absent information on inbound arrivals; Ladier and Alpan (2016b)
develop robust schedules with time windows; and Xi et al.~(2020) address
uncertain arrival and operation times through two-stage conflict robust
optimization. Our model is deterministic: each release time is known
when the schedule is constructed. These uncertainty models therefore
provide natural extensions rather than direct substitutes for the
problem studied here.

\textbf{Compound trucks and partial unloading.} Joo and Kim (2013)
introduced compound trucks with exclusive door service. Shahmardan and
Sajadieh (2020) added partial unloading, a mixed door service mode, a
MILP, a VAA-based constructive heuristic, and simulated annealing
variants whose neighborhood selection is driven by multi-armed bandits
or Q-learning; their SA-RL5 variant is our main baseline. Follow-up work
on this model line is sparse, and none considers arrival times or due
dates.

\textbf{Closest neighboring problem.} Li et al.~(2026) study integrated
truck assignment and scheduling with \emph{mixed service mode docks} and
time windows, solved by a Q-learning-based adaptive large neighborhood
search. Their flexibility lives at the dock level (a door may serve
either direction), while trucks remain purely inbound or outbound and
transfers are routed through storage by AGVs; in our problem the
flexibility lives at the truck level (compound trucks, direct
door-to-door transfer). Their Q-learning, like the original model's, is
trained online within each instance.

\textbf{Learned metaheuristics.} Neural construction policies (Kool et
al., 2019; Zhang et al., 2020) and learned operator selection (via
bandits, tabular Q (Watkins \& Dayan, 1992), or deep Q-networks (Mnih et
al., 2015)) have been reported to improve many metaheuristics. The
experiments below examine a less commonly reported outcome: after
strengthening the underlying local search, the learned selector adds
little on this benchmark.

\subsection{3. Problem definition}\label{problem-definition}

\subsubsection{3.1 Base model}\label{base-model}

Let I be the set of compound trucks, F the outbound trucks, D the
destinations (\textbar I\textbar{} + \textbar F\textbar{} =
\textbar D\textbar), K the product types, and M the dock doors
(\textbar I\textbar{} ≤ \textbar M\textbar). A compound truck i carries
f\_idk units of product k for destination d; the unit handling time is
t\_k, and moving one batch between doors m and n takes t\_mn. Truck i
needs changeover times DE\_i (docking) and DL\_i (undocking). Define the
handling time h\_id = Σ\_k f\_idk · t\_k.

Decisions: (i) each compound truck is assigned one \emph{retained}
destination and one door (at most one compound truck per door); (ii)
each outbound truck is assigned one destination and one door; (iii)
outbound trucks sharing a door are sequenced. Each destination is served
by exactly one truck (its \emph{carrier}). A compound truck unloads
everything except its retained destination's demand, transfers move
products to carrier doors, and each carrier loads its destination's
demand collected from all compound trucks.

Timing follows the evaluator semantics of Shahmardan and Sajadieh
(2020): a compound truck's unload finishes at r\_i + DE\_i + Σ\_\{d′≠d\}
h\_id′; a destination is ready at its carrier door when the last
contributing transfer arrives; loading starts at max(own unload finish,
destination ready) for compound carriers and max(door available,
destination ready, r\_f) for outbound carriers; DL is added after
loading. The makespan τ is the largest completion time.

\subsubsection{3.2 Mixed-integer
formulation}\label{mixed-integer-formulation}

We state the formulation implemented in our MILP. Binary \(x_{idm}\)
equals one when compound truck \(i\in I\) retains destination \(d\in D\)
at door \(m\in M\); \(z_{fdm}\) analogously assigns outbound truck
\(f\in F\). Define \(X_{im}=\sum_d x_{idm}\), \(Z_{fm}=\sum_d z_{fdm}\),
and \(Z_{fd}=\sum_m z_{fdm}\). Binary \(b_{fgm}\) equals one when
outbound \(f\) precedes \(g\) on door \(m\). Continuous
\(U_i,C_i,S_f,C_f\) are compound unload completion, compound completion,
outbound start, and outbound completion; \(C_{\max}\) is makespan. Let
\(H_i=\sum_d h_{id}\), \(L_{id}=\sum_{j\ne i}h_{jd}\),
\(L_d=\sum_i h_{id}\), and let \(B\) be a valid upper bound on the
scheduling horizon.

Carrier and door assignments satisfy

\[
\sum_{d,m}x_{idm}=1\quad\forall i, \qquad
\sum_{d,m}z_{fdm}=1\quad\forall f, \tag{2}
\]

\[
\sum_{i,m}x_{idm}+\sum_{f,m}z_{fdm}=1\quad\forall d, \tag{3}
\]

\[
\sum_{i,d}x_{idm}\le1\quad\forall m. \tag{4}
\]

Compound timing and incoming-transfer precedence are

\[
U_i=r_i+DE_i+H_i-\sum_{d,m}h_{id}x_{idm}\quad\forall i, \tag{5}
\]

\[
C_i\ge U_i+L_{id}+DL_i-B(1-x_{idm})\quad\forall i,d,m, \tag{6}
\]

\[
C_i\ge U_j+t_{nm}+L_{id}+DL_i-B(2-x_{idm}-X_{jn}) \tag{7}
\]

for every \(i,d,m,n\) and contributing \(j\ne i\) with
\(\sum_k f_{jdk}>0\). Thus a compound carrier waits for every source
transfer. For outbound carriers,

\[
S_f\ge r_f\quad\forall f, \tag{8}
\]

\[
C_f\ge S_f+DE_f+L_d+DL_f-B(1-Z_{fd})\quad\forall f,d, \tag{9}
\]

\[
S_f\ge U_i+t_{nm}-B(2-z_{fdm}-X_{in}) \tag{10}
\]

for every \(f,d,m,n\) and contributing \(i\) with \(\sum_k f_{idk}>0\).
An outbound truck also waits for a compound truck assigned to the same
door:

\[
S_f\ge C_i-B(2-Z_{fm}-X_{im})\quad\forall f,i,m. \tag{11}
\]

For each outbound pair \(f<g\) sharing door \(m\),

\[
b_{fgm}+b_{gfm}\ge Z_{fm}+Z_{gm}-1, \tag{12}
\]

with \(b_{fgm}\le Z_{fm}\) and \(b_{fgm}\le Z_{gm}\). The associated
non-overlap constraints are

\[
S_g\ge C_f-B(1-b_{fgm}),\qquad
S_f\ge C_g-B(1-b_{gfm}). \tag{13}
\]

Finally,

\[
C_{\max}\ge C_i\quad\forall i,\qquad
C_{\max}\ge C_f\quad\forall f. \tag{14}
\]

Equations (2)--(14) model carrier assignment, compound-door exclusivity,
transfer precedence, door capacity, and outbound sequencing. The CP-SAT
model uses the same semantics but replaces the big-M disjunctions with
enforcement literals. Fixing all assignments to a feasible schedule
reproduces the independent evaluator's completion times, providing a
model-fidelity check.

\subsubsection{3.3 Time-window extension and
objective}\label{time-window-extension-and-objective}

Each truck \(q\in I\cup F\) has release time \(r_q\ge0\) and soft due
date \(\bar d_q\). Introduce \(T_q\ge0\) with

\[
T_q\ge C_q-\bar d_q. \tag{15}
\]

The optimization objective is

\[
\min\ C_{\max}+\lambda\sum_{q\in I\cup F}T_q. \tag{16}
\]

We use \(\lambda=1\); \(r_q=0\) and \(\bar d_q=\infty\) recover the base
model. CP-SAT scales time exactly to 0.01 units and returns both an
incumbent and a proven lower bound. We call an incumbent optimal only
when the solver proves equality with its bound.

\subsubsection{3.4 Combinatorial lower
bounds}\label{combinatorial-lower-bounds}

For gap reporting when exact bounds are weak we use two valid bounds
whose maximum bounds τ: (i) a \emph{critical-chain} bound --- for each
truck, the fastest possible completion over all retained-destination
choices, relaxing all door contention; (ii) a \emph{door-area} bound ---
the total unavoidable door occupancy divided by \textbar M\textbar. A
structural tardiness bound follows by applying (i) per truck against its
due date; since both terms bound their objective components for any
feasible solution, LB\_τ + λ·LB\_T is a valid bound on (16).

\subsection{4. VAA-GILS: guided iterated local
search}\label{vaa-gils-guided-iterated-local-search}

VAA-GILS instantiates the iterated local search framework (Lourenço et
al., 2019) and is deterministic apart from the sampling inside operators
and acceptance. One run proceeds as follows.

\textbf{Construction.} The VAA heuristic of Shahmardan and Sajadieh
(2020) assigns compound trucks to retained destinations by Vogel-style
regret (Korukoğlu \& Ballı, 2011), assigns remaining destinations to
outbound trucks by completion-time priority, seeds central doors with
the first assigned trucks, and inserts remaining outbound trucks into
the door with the earliest finish time. The regret assignment and
priority ranking retain the base heuristic's processing-time logic; once
doors are assigned, all finish-time and insertion computations enforce
the release times.

\textbf{Descent.} A best-improvement local search over three move
families: relocation of any outbound truck to any door position,
compound door swaps and moves to free doors, and destination swaps
between any two trucks. Descent terminates in a local optimum of this
neighborhood. It is applied to the initial solution, to every new
incumbent, and to the final solution.

\textbf{Iterated search.} For 1,000 iterations: an operator is drawn by
the selection policy (Section 4.1); the operator generates one candidate
from the current solution; if the candidate improves the global best,
descent is applied and the incumbent is updated; acceptance follows
simulated annealing (Kirkpatrick et al., 1983) with initial temperature
\(T_0=\max\{1,0.05J(s_0)\}\) and geometric cooling
\(T\leftarrow0.995T\).

\textbf{Restart and reheating.} Let \(s^*\) be the current best solution
and \(J(s^*)\) its objective. If 30 consecutive candidate evaluations
produce no new best, the current solution is replaced by a copy of
\(s^*\) and then perturbed by three successive kicks. Each kick
independently draws one of the seven generic neighborhood operators with
equal probability and applies the sampled move to the result of the
preceding kick. No descent is performed at the restart. After the
iteration's normal cooling, the temperature is reheated according to

\[
T\leftarrow\max\{T,\ 0.02J(s^*)\}. \tag{17}
\]

Thus reheating sets a lower temperature floor relative to the best
objective; it does not reset the temperature to \(T_0\) and never lowers
the current temperature. The counters for iterations since the last best
solution and for consecutive worsening candidates are both reset to zero
after the restart.

\textbf{Operator pool.} Seven generic neighborhood moves from Shahmardan
and Sajadieh (2020) (destination swaps, door swaps, insertions; the base
model numbers its moves k1--k8 but defines no k5, so seven are
implemented) plus four \emph{guided} operators that first identify the
current bottleneck and then perform a best-of mini-search around it:
(g1) relocate the last outbound truck on the critical door to its best
door/position; (g2) swap the critical truck's destination against all
trucks; (g3)/(g4) the analogous moves anchored on the most tardy truck,
falling back to (g1)/(g2) when no truck is late. A guided move costs
tens of evaluations; a fast incremental evaluator (tens of microseconds
per evaluation) makes both descent and guided moves cheap.

\subsubsection{4.1 Operator-selection
policies}\label{operator-selection-policies}

The selection policy is a plug-in, which lets us isolate its
contribution. The engine's default is uniform random selection; the two
learned policies below are evaluated only as ablation arms (Section
6.2), not as part of the proposed method.

\begin{itemize}
\tightlist
\item
  \textbf{Uniform} (default): uniform random choice over the operator
  pool.
\item
  \textbf{Tabular Q} (the policy of SA-RL5 in Shahmardan and Sajadieh,
  2020): Q-learning (Watkins \& Dayan, 1992) over five stagnation-bin
  states, trained online within the run; ε-greedy with roulette
  exploitation, shaped reward (2 for a new incumbent, 1 for a
  non-worsening move, 0 otherwise).
\item
  \textbf{Transfer DQN}: a two-layer deep Q-network (Mnih et al., 2015)
  over a 27-dimensional scale-invariant state (instance descriptors such
  as compound-truck fraction, flow concentration, window tightness;
  search descriptors such as progress, temperature, stagnation,
  door-load imbalance; and per-operator recent success rates). It is
  trained offline on 500 runs over a \emph{training} instance pool
  (sizes S/M, all flow patterns and window levels) and applied zero-shot
  (ε = 0, no updates) to unseen test instances.
\end{itemize}

\subsubsection{4.2 Relation to the base model of Shahmardan and Sajadieh
(2020)}\label{relation-to-the-base-model-of-shahmardan-and-sajadieh-2020}

VAA-GILS shares the construction heuristic and the seven generic
neighborhoods of the base model (Shahmardan \& Sajadieh, 2020), and like
it applies exactly one operator per iteration. The differences are
structural, and Table 1 makes them explicit: Shahmardan and Sajadieh
(2020) is a \emph{pure} simulated-annealing method in which the
reinforcement-learning policy chooses which single move to apply and the
step ends at the SA accept/reject; VAA-GILS embeds that same
one-operator step inside an iterated local search, adding
bottleneck-guided operators, a best-improvement descent after every
improving move, and kick restarts with reheating.

\textbf{Table 1.} Base model (Shahmardan \& Sajadieh, 2020) versus
VAA-GILS.

{\def\LTcaptype{none} % do not increment counter
\begin{longtable}[]{@{}
  >{\raggedright\arraybackslash}p{(\linewidth - 4\tabcolsep) * \real{0.3333}}
  >{\raggedright\arraybackslash}p{(\linewidth - 4\tabcolsep) * \real{0.3333}}
  >{\raggedright\arraybackslash}p{(\linewidth - 4\tabcolsep) * \real{0.3333}}@{}}
\toprule\noalign{}
\begin{minipage}[b]{\linewidth}\raggedright
Aspect
\end{minipage} & \begin{minipage}[b]{\linewidth}\raggedright
Base SA-RL5 (Shahmardan \& Sajadieh, 2020)
\end{minipage} & \begin{minipage}[b]{\linewidth}\raggedright
VAA-GILS (ours)
\end{minipage} \\
\midrule\noalign{}
\endhead
\bottomrule\noalign{}
\endlastfoot
Overall framework & Pure simulated annealing & Iterated local search (SA
only as acceptance) \\
Operators per iteration & One & One (same) \\
Operator pool & 7 neighborhoods (numbered k1--k8, k5 unused) & Those 7 +
4 bottleneck-guided (g1--g4) \\
Guided-move mechanics & Roulette-wheel single move (k6--k8) &
Deterministic best-of scan around the makespan-critical / most-tardy
truck \\
Which operator fires & Online Q-learning / roulette & Uniform random
(learned selection kept only as ablation; no benefit) \\
After an improving move & Nothing --- the SA step is the whole iteration
& Full best-improvement descent to a local optimum \\
Stagnation handling & None & Kick restart from best + reheating \\
Acceptance & SA (Metropolis) & SA (Metropolis) --- shown dispensable by
ablation (Section 6.2) \\
Construction & VAA & VAA, extended with release times \\
Objective & Makespan & Makespan + λ · total tardiness (time windows) \\
\end{longtable}
}

The most consequential difference in Table 1 occurs after an improving
move. In VAA-GILS, that move triggers descent, so the selected operator
often begins an improvement that is completed by several further moves.
The ablation later shows that this descent accounts for most of the
quality gain. Operator choice is a separate issue: Shahmardan and
Sajadieh (2020) learn this choice, whereas our uniform control reaches
essentially the same final quality as the learned selectors.

\subsection{5. Experimental design}\label{experimental-design}

\textbf{Instances.} We follow the size regime of Shahmardan and Sajadieh
(2020): (\textbar I\textbar, \textbar D\textbar, \textbar M\textbar) = S
(6, 9, 6), M (12, 18, 12), L (20, 30, 20), with compound-truck fraction
2/3. Flows f\_idk are uniform integers in {[}0, 20{]} with 3 product
types; unit times t\_k \textasciitilde{} U{[}1, 5{]}; doors lie on a
random 100×100 layout. Window levels: \emph{none} (r = 0, d̄ = ∞),
\emph{medium}, and \emph{tight}: releases r \textasciitilde{} U{[}0,
ρH{]} and d̄ = r + δH, where H estimates the unconstrained makespan and
(ρ, δ) = (0.25, 0.60) and (0.50, 0.35) respectively.

\textbf{Protocol.} Seeds are split into disjoint train / tuning / test
pools. All algorithm design and hyperparameters were frozen on the
tuning pool; the DQN was trained on the train pool; every number below
is the single test-pool run. The main solution-quality comparison (Table
2 and its Wilcoxon tests) uses 20 instances per cell × 5 replications (9
cells: 3 sizes × 3 window levels); the controlled ablations (Tables 3--5
and Figure 3) keep the original 5-instance design, since they are
equal-budget, equal-structure decompositions rather than generalization
claims. CP-SAT runs once per instance with 8 threads: 300 s on S (5
instances) and 600 s on M and L (2 instances per cell). λ = 1 on window
cells.

\textbf{Statistics.} The independently generated instance is the
experimental unit. For stochastic methods, we first average the five
replications within each instance and method, then apply two-sided
Wilcoxon signed-rank tests to paired instance means. Zero differences
are omitted by the signed-rank test, so the reported \(n\) can be
smaller than the number of paired instances. Replications quantify
algorithmic variability but are not treated as independent samples.

\subsection{6. Results}\label{results}

\subsubsection{6.1 Solution quality}\label{solution-quality}

\emph{Scope of claims.} We report proximity to CP-SAT incumbents on the
S cells (300 s per instance) and M-none (600 s, two instances). Only the
five S-none incumbents were proven optimal, so only that subset supports
an optimality-gap statement; all other exact-solver comparisons are
incumbent gaps. On the remaining cells CP-SAT returns no feasible
solution within 600 s; there our claim is dominance over the tested
baselines, and the reported solutions are, to our knowledge, the best
known. Accordingly, Table 2 reports the gap to the \emph{best-known}
solution per instance (Δbk, over all methods, all budgets, and CP-SAT)
as the primary quality measure, and the gap to the CP-SAT reference
where one exists.

\textbf{Table 2.} Test-pool results. Mean obj ± std and Δbk are over 20
instances × 5 reps per cell (Phase D1); Δbk = mean gap to the
per-instance best-known solution (over all methods and CP-SAT where
available). vs CP-SAT is over the exact-referenced subset only (5 per S
cell --- S-none: proven optima; 2 for M-none) and is unchanged from the
original exact runs.

{\def\LTcaptype{none} % do not increment counter
\begin{longtable}[]{@{}lllrr@{}}
\toprule\noalign{}
Cell & Method & Mean obj ± std & Δbk (\%) & vs CP-SAT (\%) \\
\midrule\noalign{}
\endhead
\bottomrule\noalign{}
\endlastfoot
S-none & VAA & 1497.5 ± 400.3 & 6.60 & +6.49 \\
& Paper-SA-RL5 & 1423.3 ± 383.6 & 1.03 & +1.21 \\
& GILS-uniform & 1411.8 ± 382.7 & 0.16 & +0.21 \\
& GILS-tabular & 1412.9 ± 381.3 & 0.28 & +0.33 \\
& GILS-DQN & 1416.1 ± 382.4 & 0.51 & +0.75 \\
S-medium & VAA & 5447.5 ± 1468.9 & 3.89 & +3.33 \\
& GILS-uniform & 5255.7 ± 1413.3 & 0.08 & +0.11 \\
& GILS-tabular & 5259.1 ± 1412.3 & 0.15 & +0.18 \\
& GILS-DQN & 5270.8 ± 1415.3 & 0.38 & +0.55 \\
S-tight & VAA & 9276.8 ± 2387.6 & 2.90 & +2.17 \\
& GILS-uniform & 9010.8 ± 2206.1 & 0.09 & +0.21 \\
& GILS-tabular & 9011.9 ± 2206.7 & 0.11 & +0.17 \\
& GILS-DQN & 9022.0 ± 2205.3 & 0.23 & +0.34 \\
M-none & VAA & 3030.9 ± 903.6 & 3.73 & +4.80 \\
& Paper-SA-RL5 & 2976.4 ± 864.4 & 1.86 & +0.99 \\
& GILS-uniform & 2928.9 ± 842.6 & 0.29 & +0.13 \\
& GILS-tabular & 2932.1 ± 846.3 & 0.37 & +0.11 \\
& GILS-DQN & 2932.5 ± 845.5 & 0.40 & +0.19 \\
M-medium & VAA & 23025.5 ± 4075.4 & 2.56 & --- \\
& GILS-uniform & 22526.4 ± 3960.5 & 0.29 & --- \\
& GILS-tabular & 22515.5 ± 3964.0 & 0.24 & --- \\
& GILS-DQN & 22540.3 ± 3956.4 & 0.36 & --- \\
M-tight & VAA & 38903.7 ± 11283.5 & 1.25 & --- \\
& GILS-uniform & 38528.0 ± 11078.0 & 0.19 & --- \\
& GILS-tabular & 38512.1 ± 11059.5 & 0.16 & --- \\
& GILS-DQN & 38543.6 ± 11077.6 & 0.24 & --- \\
L-none & VAA & 5492.3 ± 1482.2 & 2.21 & --- \\
& Paper-SA-RL5 & 5479.7 ± 1449.5 & 1.97 & --- \\
& GILS-uniform & 5388.4 ± 1427.3 & 0.27 & --- \\
& GILS-tabular & 5385.5 ± 1424.6 & 0.22 & --- \\
& GILS-DQN & 5386.1 ± 1425.2 & 0.23 & --- \\
L-medium & VAA & 55611.1 ± 18087.3 & 1.52 & --- \\
& GILS-uniform & 54846.4 ± 17523.5 & 0.10 & --- \\
& GILS-tabular & 54833.4 ± 17526.8 & 0.07 & --- \\
& GILS-DQN & 54848.8 ± 17522.9 & 0.11 & --- \\
L-tight & VAA & 119671.9 ± 32475.3 & 0.77 & --- \\
& GILS-uniform & 118901.4 ± 31781.1 & 0.08 & --- \\
& GILS-tabular & 118865.6 ± 31788.4 & 0.05 & --- \\
& GILS-DQN & 118888.8 ± 31789.1 & 0.07 & --- \\
\end{longtable}
}

On every cell with an exact reference, GILS approaches it (see the final
column of Table 2): within 0.21--0.75\% of the five proven optima on
S-none, within 0.11--0.55\% of the 300 s CP-SAT incumbents on S-medium
and S-tight, and within 0.11--0.19\% of the 600 s incumbents on M-none,
which the best GILS runs match or beat. The ordering GILS \textless{}
Paper-SA-RL5 \textless{} VAA holds in every cell where the baseline is
defined; over the 20-instance test set, Wilcoxon tests on paired
instance means confirm GILS-tabular beats VAA (+2.64\% mean, n = 180, p
\textless{} 10⁻⁴) and Paper-SA-RL5 (+1.33\%, n = 60, p \textless{}
10⁻⁴). The combinatorial lower bounds are informative on no-window cells
but loose elsewhere: the gap of the \emph{best-known solution itself} to
the bound is 0.0\% on S-none (optimality proven), 18.5--51.5\% on the
other S cells, 33--34\% on M-none and L-none, and 171--273\% on the
time-window cells at M and L. These figures measure the bound, not the
solutions; tightening bounds for large time-window instances is left as
future work.

\textbf{Runtime.} Mean per-run times of GILS (1,000 iterations): 0.11 s
(S), 0.52 s (M), 2.3 s (L). CP-SAT: 76 s (S, mean incl.~early optimality
proofs), 237 s (M), 376 s (L). GILS attains its quality at roughly
70--700× less computation.

\subsubsection{6.2 Where does the quality come
from?}\label{where-does-the-quality-come-from}

All ablations use 1,000 iterations. We change one part of the engine at
a time: the available moves, the search components, or the
operator-selection rule. Uniform sampling is retained unless the
selector itself is under study. Thus, the paired comparisons differ only
in the component named in each table.

\textbf{Operator pool (guided operators).} We fix the selection policy
to uniform and vary only the operator pool: \emph{generic} (the seven
paper neighborhoods), \emph{critical} (generic plus the two
makespan-critical guided operators g1, g2), and \emph{full} (critical
plus the two tardiness-guided operators g3, g4). The mean gap to the
per-instance best-known solution is monotone --- generic ≥ critical ≥
full --- in every one of the nine cells (e.g., S-none 0.45 / 0.17 /
0.10; M-medium 0.45 / 0.37 / 0.28). Table 3 gives the paired tests.

\textbf{Table 3.} Operator-pool ablation (two-sided Wilcoxon on paired
instance replication means; +mean = the richer pool is better).

{\def\LTcaptype{none} % do not increment counter
\begin{longtable}[]{@{}
  >{\raggedright\arraybackslash}p{(\linewidth - 10\tabcolsep) * \real{0.1429}}
  >{\raggedright\arraybackslash}p{(\linewidth - 10\tabcolsep) * \real{0.1429}}
  >{\raggedleft\arraybackslash}p{(\linewidth - 10\tabcolsep) * \real{0.1905}}
  >{\raggedleft\arraybackslash}p{(\linewidth - 10\tabcolsep) * \real{0.1905}}
  >{\raggedleft\arraybackslash}p{(\linewidth - 10\tabcolsep) * \real{0.1905}}
  >{\raggedright\arraybackslash}p{(\linewidth - 10\tabcolsep) * \real{0.1429}}@{}}
\toprule\noalign{}
\begin{minipage}[b]{\linewidth}\raggedright
Contrast
\end{minipage} & \begin{minipage}[b]{\linewidth}\raggedright
Scope
\end{minipage} & \begin{minipage}[b]{\linewidth}\raggedleft
n
\end{minipage} & \begin{minipage}[b]{\linewidth}\raggedleft
Mean (\%)
\end{minipage} & \begin{minipage}[b]{\linewidth}\raggedleft
p
\end{minipage} & \begin{minipage}[b]{\linewidth}\raggedright
Verdict
\end{minipage} \\
\midrule\noalign{}
\endhead
\bottomrule\noalign{}
\endlastfoot
generic → critical (add g1, g2) & all & 40 & +0.114 & \textless10⁻⁴ &
significant \\
critical → full (add g3, g4) & all & 42 & +0.047 & 0.0001 &
significant \\
critical → full & TW cells & 28 & +0.045 & 0.0003 & significant \\
generic → full (all guided) & all & 43 & +0.161 & \textless10⁻⁴ &
significant \\
\end{longtable}
}

\emph{{[}Figure --- supplied as a separate vector file (EPS/PDF) at
submission.{]}}

\textbf{Figure 1.} Operator-pool ablation. Adding the bottleneck-guided
operators lowers the mean gap to best-known in every one of the nine
cells; the ordering generic \textgreater{} critical \textgreater{} full
is monotone throughout.

The guided operators lower the objective consistently and significantly,
and the effect never flips sign. (On no-time-window cells the g3/g4 gain
is partly a selection-weight effect, since they fall back to g1/g2 when
no truck is late; the unambiguous new-capability results are generic →
critical everywhere and critical → full on the time-window cells.)

\textbf{Engine components.} Holding the pool at \emph{full} and the
policy at uniform, we remove one engine component at a time and measure
the resulting degradation (Table 4). Best-improvement descent is the
dominant contributor and kick restarts are second, both significant
everywhere. The Vogel-approximation initial solution has no
statistically detectable pooled effect --- starting from a random
feasible solution reaches essentially the same objective --- consistent
with a separate sensitivity test in which a purely random start
(26--31\% worse than VAA) changes the final objective by at most
±0.33\%. Removing stochastic acceptance (making it greedy) does
\emph{not} worsen the objective on average, and is marginally better at
M and L; the deterministic descent, restarts, and guided operators
already reach the attractor without it.

\textbf{Table 4.} Engine-component ablation (leave-one-out; degradation
= relative objective increase when removed; Wilcoxon on paired instance
replication means).

{\def\LTcaptype{none} % do not increment counter
\begin{longtable}[]{@{}
  >{\raggedright\arraybackslash}p{(\linewidth - 8\tabcolsep) * \real{0.1667}}
  >{\raggedleft\arraybackslash}p{(\linewidth - 8\tabcolsep) * \real{0.2222}}
  >{\raggedleft\arraybackslash}p{(\linewidth - 8\tabcolsep) * \real{0.2222}}
  >{\raggedleft\arraybackslash}p{(\linewidth - 8\tabcolsep) * \real{0.2222}}
  >{\raggedright\arraybackslash}p{(\linewidth - 8\tabcolsep) * \real{0.1667}}@{}}
\toprule\noalign{}
\begin{minipage}[b]{\linewidth}\raggedright
Component removed
\end{minipage} & \begin{minipage}[b]{\linewidth}\raggedleft
n
\end{minipage} & \begin{minipage}[b]{\linewidth}\raggedleft
Degradation (\%)
\end{minipage} & \begin{minipage}[b]{\linewidth}\raggedleft
p
\end{minipage} & \begin{minipage}[b]{\linewidth}\raggedright
Verdict
\end{minipage} \\
\midrule\noalign{}
\endhead
\bottomrule\noalign{}
\endlastfoot
best-improvement descent & 45 & +0.156 & \textless10⁻⁴ & significant \\
kick restart & 43 & +0.069 & \textless10⁻⁴ & significant \\
VAA initial solution & 45 & +0.022 & 0.114 & no significant
difference \\
stochastic acceptance & 45 & −0.026 & 0.0014 & greedy slightly better \\
\end{longtable}
}

\emph{{[}Figure --- supplied as a separate vector file (EPS/PDF) at
submission.{]}}

\textbf{Figure 2.} Engine-component leave-one-out ablation.
Best-improvement descent and kick restarts are the significant
contributors; the VAA initial solution has no statistically significant
pooled effect, and removing stochastic acceptance slightly improves the
mean objective in this experiment.

\textbf{Learned operator selection.} Finally we vary the selection
policy itself, holding pool and components at \emph{full}: uniform
random, the base model's tabular Q-learning, and a transfer-trained deep
Q-network (Table 5). No learned policy beats uniform by a practically
meaningful margin. Across budgets (mean gap to best-known) uniform runs
0.47 → 0.16\%, tabular 0.56 → 0.09\%, and DQN 0.56 → 0.26\% from 50 to
3,000 iterations: at 50 iterations uniform is \emph{better} than tabular
(+0.08\%, p = 0.0055), at 3,000 tabular edges uniform (+0.07\%, p
\textless{} 10⁻⁴), and the transfer DQN never wins at any budget and is
significantly worse at 1,000 and 3,000. All selection effects are ≤ 0.17
percentage points --- below the guided-operator effect of Table 3 and an
order of magnitude below the method-level differences of Table 2 --- and
their direction flips with budget.

Crucially, the same picture holds on the no-time-window cells alone,
which are exactly the original problem of Shahmardan and Sajadieh (2020)
with all trucks available at time zero: at 1,000 iterations tabular
Q-learning is indistinguishable from uniform (mean −0.05\%, p = 0.52, n
= 58) and the transfer DQN is significantly worse (uniform better by
+0.14\%, p = 0.0012; tabular better by +0.09\%, p \textless{} 10⁻³). The
negative result is therefore not an artifact of the time-window
extension; it already holds in the base setting where the original model
introduces learned selection.

\textbf{Table 5.} Selection-policy effects at 1,000 iterations
(two-sided Wilcoxon on paired instance replication means; positive mean
= first method better).

{\def\LTcaptype{none} % do not increment counter
\begin{longtable}[]{@{}lrrl@{}}
\toprule\noalign{}
Comparison & Mean diff (\%) & p & Verdict \\
\midrule\noalign{}
\endhead
\bottomrule\noalign{}
\endlastfoot
tabular vs uniform & −0.01 & 0.150 & no difference \\
tabular vs DQN & +0.10 & \textless10⁻⁴ & DQN worse \\
uniform vs DQN & +0.11 & \textless10⁻⁴ & DQN worse \\
\end{longtable}
}

\emph{{[}Figure --- supplied as a separate vector file (EPS/PDF) at
submission.{]}}

\textbf{Figure 3.} Selection-policy budget sweep. Across budgets from 50
to 3,000 iterations, neither learned policy beats uniform random
selection by a practically meaningful margin; the transfer DQN is
consistently worse and the tabular policy only edges uniform at the
largest budget.

Across the three experiments, descent has the largest measured effect.
Guided moves and restarts provide smaller but repeatable gains. The
initial solution, stochastic acceptance, and learned selection have
little practical effect once those parts are present. One tuning
instance illustrates the pattern: several selector variants converged to
exactly the CP-SAT incumbent obtained after 240 seconds. This is only a
point check, but it shows how the local search can erase differences
created earlier in a run.

\subsubsection{6.3 Tardiness-weight
sensitivity}\label{tardiness-weight-sensitivity}

The objective (16) weights total tardiness by λ; our experiments fix λ =
1. To justify this modelling choice we trace the makespan--tardiness
trade-off by sweeping λ ∈ \{0, 0.25, 0.5, 1, 2, 4, 8\} on the
\emph{tuning} pool (all six time-window cells, five instances × three
replications, 1,000 iterations, seeds shared across λ so that only the
objective weight differs). For each run we record the \emph{physical}
makespan and total tardiness of the returned schedule; to combine cells
of very different magnitude we min--max normalize each metric across the
λ grid per instance and average (Figure 4).

\emph{{[}Figure --- supplied as a separate vector file (EPS/PDF) at
submission.{]}}

\textbf{Figure 4.} Tardiness-weight sensitivity. As λ grows the search
trades a higher makespan for lower tardiness along a convex frontier
(normalized makespan 0.13 → 0.74, normalized tardiness 0.99 → 0.04).
Most of the achievable tardiness reduction is bought by the first
increment away from λ = 0, after which returns diminish while the
makespan penalty keeps rising.

The absolute trade-off is mild. Across the sweep, mean makespan rises by
only about 1\% (e.g., S 2060.8 → 2084.1), while mean tardiness falls by
2--3\% (e.g., L 87215 → 86306). At λ = 1, tardiness is already near its
minimum (normalized 0.14, versus 0.04 at the extreme λ = 8) at only a
moderate makespan cost (normalized 0.41). We use λ = 1 because it also
has a direct interpretation: one unit of makespan time is exchanged for
one unit of tardiness time. The sweep is reported as design
justification, not as a tuned result, and uses only the tuning pool.

\subsection{7. Discussion}\label{discussion}

The selector result should be interpreted within the search process used
here. A candidate can be evaluated in tens of microseconds, so an
unhelpful operator choice is inexpensive. More importantly, every new
incumbent is followed by a deterministic descent. Two runs that enter
the same basin can therefore end at the same schedule even if their
preceding moves differ. The budget curves suggest that this happens
before 1,000 iterations for most test instances.

There is also little information for a policy to predict. Instances are
static, and all releases are known before optimization begins. This is
quite different from online dock control, where arrivals are revealed
over time, or from a simulation model in which evaluating one move is
costly. Learned selection may matter more in either setting. It may also
matter on instances large enough that descent cannot reach saturation
within the available budget. We did not test those cases, so they remain
hypotheses rather than consequences of the present experiments.

What the results do support is a narrower recommendation. A learned
selector should be compared with uniform sampling inside the same engine
and under the same evaluation budget. Without that control, gains due to
descent, restart, or a stronger move set can be attributed to the policy
itself.

\subsection{8. Conclusion}\label{conclusion}

This study adds release times and soft due dates to compound-truck
scheduling with partial unloading. The MILP and CP-SAT formulations also
provide a check on the heuristic results. On the exact-referenced
subsets, VAA-GILS finished within a fraction of a percent of the CP-SAT
incumbent; optimality was proven only for the small no-window cases. For
the larger windowed instances, it found the best schedules among the
methods tested in a few seconds, while CP-SAT did not return a feasible
schedule within its limit.

The ablations change how that result should be understood. Most of the
gain is due to best-improvement descent, with additional gains from
bottleneck moves and restarts. In this engine, tabular Q-learning and
the transfer DQN did not improve meaningfully on uniform operator
sampling. Nor did the final result depend much on the starting solution
or simulated-annealing acceptance. These findings are specific to the
present benchmark. An online version with unrevealed arrivals would
provide a more demanding test of adaptive policies; tighter bounds for
the larger time-window cases are another priority.

\subsection{Declaration of generative AI and AI-assisted technologies in
the manuscript preparation
process}\label{declaration-of-generative-ai-and-ai-assisted-technologies-in-the-manuscript-preparation-process}

During the preparation of this work the authors used Claude Code
(Anthropic) in order to assist with software for the computational
experiments, generation of result figures, and drafting and language
editing of the manuscript. After using this tool, the authors reviewed
and edited the content as needed and take full responsibility for the
content of the publication.

\subsection{References}\label{references}

Assadi, M. T., \& Bagheri, M. (2016). Differential evolution and
population-based simulated annealing for truck scheduling problem in
multiple door cross-docking systems. \emph{Computers \& Industrial
Engineering, 96}, 149--161. https://doi.org/10.1016/j.cie.2016.03.021

Bodnar, P., de Koster, R. B. M., \& Azadeh, K. (2017). Scheduling trucks
in a cross-dock with mixed service mode dock doors. \emph{Transportation
Science, 51}(1), 112--131. https://doi.org/10.1287/trsc.2015.0612

Boysen, N., \& Fliedner, M. (2010). Cross dock scheduling:
Classification, literature review and research agenda. \emph{Omega,
38}(6), 413--422.

Joo, C. M., \& Kim, B. S. (2013). Scheduling compound trucks in
multi-door cross-docking terminals. \emph{International Journal of
Advanced Manufacturing Technology, 64}, 977--988.

Kirkpatrick, S., Gelatt, C. D., \& Vecchi, M. P. (1983). Optimization by
simulated annealing. \emph{Science, 220}(4598), 671--680.

Konur, D., \& Golias, M. M. (2013a). Analysis of different approaches to
cross-dock truck scheduling with truck arrival time uncertainty.
\emph{Computers \& Industrial Engineering, 65}(4), 663--672.
https://doi.org/10.1016/j.cie.2013.05.009

Konur, D., \& Golias, M. M. (2013b). Cost-stable truck scheduling at a
cross-dock facility with unknown truck arrivals: A meta-heuristic
approach. \emph{Transportation Research Part E: Logistics and
Transportation Review, 49}(1), 71--91.
https://doi.org/10.1016/j.tre.2012.06.007

Kool, W., van Hoof, H., \& Welling, M. (2019). Attention, learn to solve
routing problems! In \emph{Proceedings of the International Conference
on Learning Representations}.

Korukoğlu, S., \& Ballı, S. (2011). An improved Vogel's approximation
method for the transportation problem. \emph{Mathematical and
Computational Applications, 16}(2), 370--381.

Ladier, A.-L., \& Alpan, G. (2016a). Cross-docking operations: Current
research versus industry practice. \emph{Omega, 62}, 145--162.

Ladier, A.-L., \& Alpan, G. (2016b). Robust cross-dock scheduling with
time windows. \emph{Computers \& Industrial Engineering, 99}, 16--28.
https://doi.org/10.1016/j.cie.2016.07.003

Larbi, R., Alpan, G., Baptiste, P., \& Penz, B. (2011). Scheduling cross
docking operations under full, partial and no information on inbound
arrivals. \emph{Computers \& Operations Research, 38}(6), 889--900.
https://doi.org/10.1016/j.cor.2010.10.003

Li, Y., Mohammadi, M., Zhang, X., Lan, Y., \& van Jaarsveld, W. (2026).
Integrated trucks assignment and scheduling problem with mixed service
mode docks: A Q-learning based adaptive large neighborhood search
algorithm. \emph{European Journal of Operational Research, 333}(1),
117--137. https://doi.org/10.1016/j.ejor.2025.12.036

Lourenço, H. R., Martin, O. C., \& Stützle, T. (2019). Iterated local
search: Framework and applications. In \emph{Handbook of metaheuristics}
(3rd ed.). Springer.

Mnih, V., Kavukcuoglu, K., Silver, D., Rusu, A. A., Veness, J.,
Bellemare, M. G., Graves, A., Riedmiller, M., Fidjeland, A. K.,
Ostrovski, G., Petersen, S., Beattie, C., Sadik, A., Antonoglou, I.,
King, H., Kumaran, D., Wierstra, D., Legg, S., \& Hassabis, D. (2015).
Human-level control through deep reinforcement learning. \emph{Nature,
518}, 529--533.

Molavi, D., Shahmardan, A., \& Sajadieh, M. S. (2018). Truck scheduling
in a cross docking systems with fixed due dates and shipment sorting.
\emph{Computers \& Industrial Engineering, 117}, 29--40.
https://doi.org/10.1016/j.cie.2018.01.009

Perron, L., \& Didier, F. (2024). \emph{Google OR-Tools CP-SAT solver}
(Version 9.x) {[}Computer software{]}.
https://developers.google.com/optimization

Rijal, A., Bijvank, M., \& de Koster, R. (2019). Integrated scheduling
and assignment of trucks at unit-load cross-dock terminals with mixed
service mode dock doors. \emph{European Journal of Operational Research,
278}(3), 752--771. https://doi.org/10.1016/j.ejor.2019.04.028

Shahmardan, A., \& Sajadieh, M. S. (2020). Truck scheduling in a
multi-door cross-docking center with partial unloading --- Reinforcement
learning-based simulated annealing approaches. \emph{Computers \&
Industrial Engineering, 139}, Article 106134.

Van Belle, J., Valckenaers, P., \& Cattrysse, D. (2012). Cross-docking:
State of the art. \emph{Omega, 40}(6), 827--846.

Van Belle, J., Valckenaers, P., Vanden Berghe, G., \& Cattrysse, D.
(2013). A tabu search approach to the truck scheduling problem with
multiple docks and time windows. \emph{Computers \& Industrial
Engineering, 66}(4), 818--826. https://doi.org/10.1016/j.cie.2013.09.024

Watkins, C. J. C. H., \& Dayan, P. (1992). Q-learning. \emph{Machine
Learning, 8}, 279--292.

Xi, X., Liu, C., Wang, Y., \& Lee, L. H. (2020). Two-stage conflict
robust optimization models for cross-dock truck scheduling problem under
uncertainty. \emph{Transportation Research Part E: Logistics and
Transportation Review, 144}, Article 102123.
https://doi.org/10.1016/j.tre.2020.102123

Zhang, C., Song, W., Cao, Z., Zhang, J., Tan, P. S., \& Chi, X. (2020).
Learning to dispatch for job shop scheduling via deep reinforcement
learning. In \emph{Advances in Neural Information Processing Systems}
(Vol. 33).

\end{document}
